%%%%%%%%%%%%%%%%%%%%%%%%%%%%%%%%%%%%%%%%%%%%%%%%%%%%%%%%%%%%%%%%%%%%%%%%%%%
%
% Generic template for TFC/TFM/TFG/Tesis
%
% By:
%  + Javier Macías-Guarasa.
%    Departamento de Electrónica
%    Universidad de Alcalá
%  + Roberto Barra-Chicote.
%    Departamento de Ingeniería Electrónica
%    Universidad Politécnica de Madrid
%
% By: + Javier Macías-Guarasa. Departamento de Electrónica Universidad de Alcalá + Roberto Barra-Chicote. Departamento de Ingeniería Electrónica Universidad Politécnica de Madrid
%
% Based on original sources by Roberto Barra, Manuel Ocaña, Jesús Nuevo, Pedro Revenga, Fernando Herránz and Noelia Hernández. Thanks a lot to all of them, and to the many anonymous contributors found (thanks to google) that provided help in setting all this up.
%
% See also the additionalContributors.txt file to check the name of additional contributors to this work.
%
% If you think you can add pieces of relevant/useful examples, improvements, please contact us at (macias@depeca.uah.es)
%
% You can freely use this template and please contribute with comments or suggestions!!!
%
%%%%%%%%%%%%%%%%%%%%%%%%%%%%%%%%%%%%%%%%%%%%%%%%%%%%%%%%%%%%%%%%%%%%%%%%%%%

\chapter{Primeros pasos}
\label{cha:primeros-pasos}

En este capítulo te acompañamos desde la preparación del entorno hasta la generación del primer PDF. Para empezar no necesitas conocer el \texttt{Makefile} ni utilizar una terminal: puedes trabajar con tu editor habitual de \LaTeX{}.

\section{Prerrequisitos}
\label{sec:prerrequisitos}

Para compilar la plantilla necesitas una distribución de \LaTeX{} suficientemente completa y actualizada. TeX Live es una opción habitual en GNU/Linux, MiKTeX y TeX Live están disponibles para Windows, y MacTeX proporciona una distribución adaptada a macOS. Asegúrate de que la distribución que elijas proporcione \texttt{pdflatex} y Biber, además de los paquetes cargados desde \texttt{Config/preamble.tex}.

Biber es necesario cuando el documento incluye bibliografía. Si utilizas acrónimos o símbolos, también necesitarás \texttt{makeglossaries}; el libro de ejemplo distribuido con la plantilla contiene ambos elementos y, por tanto, requiere esta herramienta para generar todos sus listados.

Si durante la compilación recibes un aviso de que falta un fichero \texttt{.sty}, instala el paquete de tu distribución de \LaTeX{} que lo contiene. En MiKTeX puedes activar la instalación automática de paquetes; TeX Live y MacTeX proporcionan sus propias herramientas para administrar los paquetes instalados.

No necesitas utilizar \texttt{make} para trabajar con la plantilla. Puedes compilarla con cualquier editor o herramienta de construcción de \LaTeX{} que permita configurar las herramientas descritas en la sección~\ref{sec:compilacion}; si lo prefieres, también puedes hacerlo desde una terminal. El \texttt{Makefile} incluido es únicamente una forma opcional de automatizar el proceso y resulta especialmente cómodo en sistemas GNU/Linux.


\section{Compilación}
\label{sec:compilacion}

Antes de compilar el documento, completa la configuración compartida en \texttt{Config/myconfig.tex} siguiendo las indicaciones de la sección~\ref{sec:definicion-del-tipo}. Después, abre \texttt{Book/book.tex} como documento principal en tu editor de \LaTeX{} y utiliza su función habitual de compilación o construcción.

Para generar correctamente todos los elementos del libro, configura el editor para utilizar \texttt{pdflatex} como compilador y Biber como procesador de la bibliografía. Si utilizas acrónimos o símbolos, añade también \texttt{makeglossaries} a la cadena de construcción. Finalmente, realiza las pasadas adicionales de \texttt{pdflatex} necesarias para actualizar las referencias cruzadas, los índices, la bibliografía y los glosarios.

\subsection{GNU/Linux}

En GNU/Linux puedes utilizar tu editor de \LaTeX{} preferido. Asegúrate de abrir \texttt{Book/book.tex} como documento principal y de configurar la cadena de construcción como se explica en la sección~\ref{sec:compilacion}.

Si prefieres automatizar el proceso mediante \texttt{make}, consulta \textit{Compilación mediante el \texttt{Makefile}} (apartado~\ref{sec:compilacion-makefile}), donde se describen el procedimiento y los objetivos disponibles.

\subsection{Windows}

En Windows puedes utilizar MiKTeX o TeX Live junto con TeXstudio u otro editor de \LaTeX{}. En la configuración del editor, selecciona Biber como procesador de la bibliografía y, si utilizas acrónimos o símbolos, añade \texttt{makeglossaries} a la cadena de construcción. Para ejecutar \texttt{makeglossaries} también necesitas un intérprete de Perl; Strawberry Perl es una de las alternativas disponibles para Windows.

No necesitas instalar \texttt{make} para trabajar con la plantilla en Windows. Si quieres utilizar el \texttt{Makefile} descrito en el apartado~\ref{sec:compilacion-makefile}, puedes hacerlo desde un entorno compatible que proporcione las herramientas requeridas, como WSL, aunque normalmente te resultará más sencillo configurar la cadena de compilación en el propio editor.

Parte de estas instrucciones se basa en la información que nos proporcionó Cristina Losada sobre la configuración de TeXmaker en Windows. Muchas gracias por su ayuda.

\subsection{macOS}

En macOS puedes utilizar MacTeX junto con TeXShop, TeXstudio u otro editor compatible. Configura la cadena de construcción para incluir \texttt{pdflatex}, Biber y, si utilizas acrónimos o símbolos, \texttt{makeglossaries} en el orden explicado en la sección~\ref{sec:compilacion}.

Si tu editor no actualiza correctamente los acrónimos o símbolos, utiliza desde el propio editor la opción para ejecutar \texttt{makeglossaries} y realiza después dos nuevas pasadas de \texttt{pdflatex}. No necesitas borrar el resto de los ficheros del proyecto para regenerar los glosarios.

Las primeras comprobaciones de la generación de glosarios en macOS fueron realizadas y compartidas con nosotros por Carlos Cruz. Muchas gracias por su ayuda.


%%% Local Variables:
%%% TeX-master: "../book"
%%% End:
